% EN49 c00 | 原书 PDF 1–11（原书 i–ix） | 前置页：摘要、致谢（封面/版权/目录由 LaTeX 生成）
% 骨架文件：标题行已设定，正文由译者填写。每页起首用 % ===== 原书 p.N ===== 注释标记。
\section*{摘要}
\addcontentsline{toc}{section}{摘要}

% ===== 原书 p.iii（PDF 5） =====
本论文就两类不同的超曲面，证明了两个同调镜像对称结果。在第一部分，我们考察 Auroux 提
出的奇异超曲面深谷范畴（Fukaya category）的一种定义：把邻近纤维的深谷范畴在 Seidel 自
然变换处作范畴局部化；并证明这一定义具有若干良好性质。首先，我们证明 Orlov 导出
Kn{\"o}rrer 周期性定理的 A 侧类比，即证明 Auroux 的范畴与某个高维 Landau--Ginzburg
模型的 Fukaya--Seidel 范畴导出等价。其次，我们说明这一定义如何蕴含某些阿贝尔簇在大复
结构极限处退化的同调镜像对称。

% ===== 原书 p.iv（PDF 6） =====
第二部分（与 Benjamin Gammage 的合作工作）证明：极仿射超曲面的同调镜像对称尊重某些自
然的辛操作，这里把镜像对称看作部分缠绕深谷范畴之间的函子；由此验证 Auroux 的若干猜想。
这些猜想关乎极仿射超曲面的镜像对称与其补集的镜像对称之间的相容性，而补集本身也是一个
极仿射超曲面。我们发现，极仿射超曲面的补集其实有两个自然镜像，其中之一是导出概型。这
两个镜像由一个以 Kn{\"o}rrer 周期性为中介的非几何等价相联系；要把这一点纳入考虑，
Auroux 的猜想就需要作出若干修正。我们的证明还用到粘合 Liouville 扇形的若干新辛技巧，
它们或许有独立的价值。

% \transnote 汇总（c00 摘要部分）
\transnote{易错预警：原文 “hypersurfaces of different kinds” 指两类超曲面（奇异超曲面、
极仿射超曲面），而不是“若干种不同维数的超曲面”。}
\transnote{含糊点拨：本译文中“深谷范畴”对应 Fukaya category，取日文汉字人名“深谷”的
通行写法；作者在正文中对它与 wrapped Fukaya 范畴并不严格区分，记号统一为
$\cW$，读者须依上下文判断。}
\transnote{背景直觉：A 侧指辛几何一侧，B 侧指代数几何一侧；“A 侧类比”意为把原本陈述
在导出范畴一侧（B 侧）的定理，改述为深谷范畴一侧（A 侧）的相应命题。}
\transnote{易错预警：“导出等价”与后文的“拟等价”不是同一个概念。前者指导出范畴之间的
等价，后者是 $A_\infty$ 范畴之间的弱等价；作者在本论文中对二者用法有别。}
\transnote{含糊点拨：“在 Seidel 自然变换处作范畴局部化”一句压缩了 §1.2 给出的定义。
直观上，这条自然变换记录了绕奇异纤维一周的单值变换；把可逆化施加其上，相当于宣告
“绕行一周不产生新信息”，从而抹去退化方向。}
\transnote{背景直觉：Landau--Ginzburg 模型是带全纯函数的代数簇，其镜像对应 Fano 簇。
作者把它当作奇异超曲面的“平滑替身”：奇异纤维难处理，而与它同伦的 LG 模型 $X\times\C$
是光滑的，故可在后者上作计算。}
\transnote{易错预警：“大复结构极限”是复结构模空间边界上的一类退化，缩写常作 LCSL；
不要与“大体积极限”（large volume limit）混淆，后者属 K{\"a}hler 模空间一侧。}
\transnote{背景直觉：极仿射超曲面指 $(\C^*)^n$ 中的代数超曲面。取补集后维数升高一维，
却仍是同一类对象，这正是 §1.1 所述“补集也有镜像”的结构性来源。}
\transnote{易错预警：原文说补集“有两个自然镜像，其中之一是导出概型”，意思是两个镜像
之中有一个不是概型而是导出概型；不要读成“补集有一个导出概型的镜像”这种单数说法。}
\transnote{含糊点拨：“非几何等价”指该等价不能由两侧几何空间之间的映射实现，只是范畴
层面的等价；它由 Orlov 的导出 Kn{\"o}rrer 周期性给出，故称“以 Kn{\"o}rrer 周期性为
中介”。}
\transnote{背景直觉：Liouville 扇形是带角流形的辛块，可像积木一样粘合；把局部模型粘成
整体是本文第二部分的证明策略，作者称这些技巧“有独立价值”，即不依赖镜像对称也成立。}

\section*{致谢}
\addcontentsline{toc}{section}{致谢}

% ===== 原书 p.viii（PDF 10） =====
首先，我要感谢导师 Denis Auroux：过去六年里，他耐心始终不渝，善意从未间断，指导无比
珍贵；他还通读了我论文中几篇文章的许多稿本。当初他愿意为一个来自新西兰的陌生学生出面
奔走，让我先到 Berkeley，再到 Harvard；若没有这份支持，我今天不会写下这些文字。我不
可能期待更好的导师。

我还要感谢合作者 Ben Gammage，他总愿意迁就我那些傻气的想法。以下各位（排名不分先后）
给了我许多有益的建议、讨论、合作与通信，谨此致谢：David Nadler、Daniel Chupin、Kevin
Lin、Sheel Ganatra、Andrew Hanlon、Mohammed Abouzaid、David Favero、Daniel
{\'A}lvarez-Gavela、Paul Hacking、John Pardon、Vivek Shende、Zack Sylvan、Jeff Hicks、
Paul Seidel、Shaoyun Bai、Abigail Ward、Yuan Yao、Ziwen Zhao。我还要感谢一位匿名审稿
人，他为本论文第一部分的文章提出了许多有益的评论与建议。

% ===== 原书 p.ix（PDF 11） =====
本论文的完成离不开新西兰皇家学会 Rutherford 基金会的支持；我要特别感谢其项目经理
Troels Petersen 的种种帮助。本工作还部分获得 NSF 资助 DMS-1937869，以及经 Denis
Auroux 获得的 Simons 基金会资助 \#385573 的支持。

% \transnote 汇总（c00 致谢部分）
\transnote{含糊点拨：原文 “reading many many drafts” 表明导师逐稿细读；“drafts of the
papers in this thesis” 指本论文所含各篇文章的稿本，不是本论文本身的稿本。}
\transnote{背景直觉：“go to bat for” 是棒球习语，意为替人出头、为其力争；此处指 Auroux
为作者争取到 Berkeley 与 Harvard 的求学机会，故译作“出面奔走”。}
\transnote{易错预警：Rutherford Foundation 是新西兰皇家学会（Royal Society of New
Zealand）下属基金会，不是英国 Rutherford 实验室或卢瑟福基金；专名一律保留英文。}
